\clearpage
\fancyhead[L]{\small CARLINI. TAVOLE ASTRONOMICHE.}
\hypertarget{native-548-lower}{}\label{native-548-lower}
\begin{center}\small Göttingische gelehrte Anzeigen. Stück 193. Seite 1921..1928. 1812 December 3.\\[3pt]\rule{.72\linewidth}{.3pt}\end{center}

\emph{Esposizione di un nuovo metodo di costruire le tavole astronomiche applicuto alle tavole del sole di} \textsc{Francesco Carlini}. 1810. \emph{Dalla reale stamperia.} LXI und 92 Seiten klein Quart. Mailand.

Durch die Arbeiten \textsc{Laplace}'s, \textsc{Delambre}'s und \textsc{von Zach}'s haben unsere neuesten Sonnentafeln einen solchen Grad von Genauigkeit erhalten, dass in dieser Hinsicht wenig mehr zu wünschen übrig bleibt. Wohl aber war für das so oft auszuführende und dem rechnenden Astronomen so viele Zeit kostende Geschäft, für die Berechnung von Sonnenörtern, noch mehr \emph{Bequemlichkeit} zu wünschen. Bei \textsc{von Zach}'s Tafeln war es weniger auf diese, als vielmehr darauf angesehen, die Tafeln in einen kleinen Raum zusammen zu drängen. Etwas mehr hat \textsc{Delambre} für den bequemen Gebrauch gesorgt, indem er die mittlern Bewegungen für alle Tage des Jahres, eine sehr detaillirte Tafel für die Mittelpunktsgleichung, und Tafeln zu doppelten Eingängen für die Störungen durch die andern Planeten lieferte. Dessen ungeachtet lassen diese Tafeln noch viel zu wünschen übrig; auch ohne Hauptänderungen zu erleiden, hätten sie an Bequemlichkeit gewonnen, wenn es ihrem Verfasser beliebt hätte, die Störungen durch den Mond vorher umzuschmelzen; mehrere kleine Tafeln nicht durch eine sehr unzeitige Sparsamkeit von ihrem wahren Platze zu versetzen, wodurch man jedesmal zu einem lästigen Hin- und Herblättern genöthigt wird; die Störungen des Radius vector zu Störungen seines Logarithmen umzuformen. Dazu kommt noch, dass die Störungen des Radius vector durch den Jupiter unrichtig und nach einer unvollständigen Formel construirt sind, wodurch man genöthigt wird, in sein Exemplar erst noch eine Correctionstafel nachzutragen.

\SourceLeaf{549}Wir freuen uns, in der vorliegenden Arbeit \textsc{Carlini}'s den Astronomen neue Sonnentafeln ankündigen zu können, die in Hinsicht auf Bequemlichkeit alles übertreffen, was man bisher gehabt, ja selbst alles, was man nur gewünscht hat. Man kann sagen, dass ihr Verfasser mit raffinirter Kunst alles gethan hat, was zur Erleichterung des Gebrauchs nur möglich war: fast jedem Wunsche ist zuvor gekommen. Die Berechnung eines Sonnenortes ist dadurch zu einem leichten Spiel geworden, wobei man kaum eine Ziffer mehr zu schreiben hat, als unumgänglich nöthig ist. Von einem solchen Werke, wofür die rechnenden Astronomen dem Verfasser den grössten Dank schuldig sind, können wir nicht umhin, eine etwas ausführlichere Anzeige zu geben.

Die Hauptidee, wovon \textsc{Carlini} einen so glücklichen Gebrauch gemacht hat, besteht darin, die Winkelargumente, wovon eine Tafel abhängt, nicht nach Graden oder nach Tausendtheilchen des Umfanges, wie es sonst üblich war, sondern nach \emph{Tagen} fortschreiten zu lassen. So lässt sich die Mittelpunktsgleichung der Sonne für eine bestimmte Excentricität (\textsc{Carlini} hat die für 1810 gewählt) in einer Tafel angeben, deren Argument die Zahl der seit dem letzten Durchgange durchs Perigeum verflossenen Tage ist; wir wollen, da von dem Zeitpunkte eines solchen Durchganges noch öfters die Rede sein muss, ihn Kürze halber mit $T$ bezeichnen. Dieses Argument wird sich ohne alle Mühe bilden, wenn man die Anzahl von Tagen kennt, die zu Anfange jeden Jahres (\textsc{Carlini} nimmt immer den Mittag des zunächst vorhergegangenen 31. December, Mailänder Meridian, dafür, ohne die Schaltjahre auszunehmen) seit dem letzten $T$ verflossen waren, oder auch bis zum nächsten $T$ noch verfliessen mussten. Für jedes Jahr kann dieses Zeitintervall, welches im ersten Falle als positiv, im zweiten als negativ betrachtet werden kann, als eine Art Epacte in einer besondern Tafel angesetzt werden, die, weil zufälliger Weise $T$ immer nahe beim Anfange des Jahres fällt, durch eine kleine Zahl vorgestellt wird. \textsc{Carlini} hat sie, um sie immer positiv zu machen, um zwei Einheiten vermehrt, und so in der Tafel II. unter der Ueberschrift A. angesetzt. Also $g$ Tage nach Anfang eines Jahres sind $g+A-2$ Tage nach dem nächsten $T$, und die dazu gehörige Mittelpunktsgleichung, sofort vereinigt mit der mittlern Bewegung während dieser $g+A-2$ Tage, findet man neben dem Argument $g+A$ in der Tafel III., unter der Ueberschrift: Prima equazione della longitudine del Sole. Hierzu braucht man also nur die Länge des Perigeum zur Zeit $T$, welche sich auf den Anfang des vorgegebenen Jahres bezieht, und die in der Tafel II. als Termine constante della longitudine del Sole mit vorkommt, zu addiren, um sofort die elliptische Länge der Sonne zu erhalten. Das eben Gesagte würde die vollständige Erklärung der beiden zuletzt genannten Rubriken in Tafel III. und II. sein, wenn die Bewegung der Sonne keine andere Ungleichheiten hätte, und die Excentricität unveränderlich wäre; allein da sich dies nicht so verhält, leiden dieselben noch einige Modificationen, welche wir jetzt erklären wollen. Die Säcular-Aenderung der Excentricität macht eine Verbesserung der Mittelpunktsgleichung nothwendig, welche durch die Formel $(t+f-1810)c$, vorgestellt werden kann, wo $t$ die Jahrzahl, $f$ der Bruch des Jahrs, $c$ eine von der mittlern Anomalie abhängige Grösse ist, welche zwischen den Grenzen $+0''173$ und $-0''173$ liegt, so dass $c = c^1 - 0''173$ gesetzt, $c^1$ immer positiv wird, und die ganze Correction \[ = (t-1810)c^1 - (t-1810)0''173 + fc \] der Theil $-(t-1810)0''173$ ist sogleich mit der Länge des Perigeum vereinigt in dem Termine constante der Tafel II. mit enthalten; der Theil $fc$ ist mit der Prima equazione der Tafel III. verbunden; auf diese Weise braucht also nur noch $(t-1810)c^1$ hinzugefügt zu werden, wozu man das immer positive $c^1$ in der Tafel IV. unter der Ueberschrift: Variazione annua della prima equazione, mit dem Argument der dritten Tafel $g+A$ findet. Ueberdies ist mit der Prima equazione in Tafel III. noch derjenige Theil der Nutation vereinigt, welcher von der Länge der Sonne abhängt. Die Tafel selbst geht \SourceLeaf{550}von Zehntheil zu Zehntheil des Tages bis 370 und verstattet also, da die erste Differenz nicht über $367''$ geht, und die zweite ganz verschwindet, eine sehr bequeme Interpolation.

Mit Erklärung der beiden Rubriken der Tafel II. sind wir noch nicht ganz fertig. Sehr sinnreich ist die Art, wie \textsc{Carlini} die Berechnung der Aberration mit Rücksicht auf ihren veränderlichen Theil ausführt. Die bisherigen Tafeln sind so eingerichtet, dass die berechnete Sonnenlänge den constanten Theil dieser Aberration $-20''25$ mit einschliesst; in ihnen ist also eigentlich \emph{sowohl} die Epoche der mittlern Länge, als die Länge des Perigeum schon um $20''25$ vermindert, und man muss dann, wenn man anders die scheinbare Länge der Sonne sucht, noch den veränderlichen Theil der Aberration aus einer besondern Tafel hinzusetzen. Diese letztere Operation hätte man, nach \textsc{Carlini}'s Bemerkung, erspart, wenn man die Länge des Perigeum nur um die Hälfte, nemlich $10''12$, vermindert hätte, indem man die um $20''25$ verminderte Länge der mittlern Länge beibehielt. Deswegen musste \textsc{Carlini} die Zahlen in der Columne der II. Tafel, die Termine constante überschrieben sind, um $10''1$ vermehren und die Zahlen $A$ derselben Tafel um 0,00286 Tag vermindern. So erhält man ohne Weiteres die mit der Aberration afficirte Sonnenlänge. Verlangt man hingegen die wahre Sonnenlänge, so hat man nichts weiter zu thun, als $10''1$ zum Termine constante, und 0,00286 Tag zu $A$ zu addiren. Endlich schliesst der Termine constante noch die Grösse $-56''6$ ein, um die Summa aller Constanten, wodurch \textsc{Carlini} die sogleich näher zu betrachtenden Störungsgleichungen positiv macht, und die $+56''6$ beträgt, wieder gut zu machen.

Auf ähnlichen Gründen, wie die Tafel für die Mittelpunktsgleichung, beruhen auch die Tafeln für die kleinen Ungleichheiten der Sonnenlänge. Die vom Monde herrührenden sind durch eine zweckmässige Verwandlung auf Gleichungen mit einfachen Argumenten zurück geführt, von welchen \textsc{Carlini} die beiden grössten, im Maximum $7''5$ und $0''5$ betragenden, beibehalten hat. In diesen Tafeln sowohl, als in der Tafel für die Nutation, ist das Argument in Tagen ausgedrückt, welche also eigentlich die Zwischenzeit bedeuten zwischen dem Tage, für welchen man rechnet, und demjenigen, wo das Argument jeder Tafel 0 war. Wie viel jedes so ausgedrückte Argument zu Anfang jeden Jahres, von 1750 bis 1900 beträgt, ist gleich in der Tafel II. mit angegeben.

Ganz vorzüglich bequem sind endlich die Tafeln für die Störungen durch die Planeten. \textsc{Carlini} hat sie zu doppelten Eingängen berechnet, welche bekanntlich die mittlere Länge der Erde und die mittlere heliocentrische Elongation des Planeten von der Erde bedeuten. Beide lassen sich wiederum durch Tage ausdrücken: aber \textsc{Carlini} hat die Tafel so angeordnet, dass man statt des ersten Arguments nur mit der Zahl von Tagen, die seit dem Anfange des Jahrs verflossen sind, und statt des zweiten mit demjenigen Werthe einzugehen hat, den eben dieses zu Anfang des Jahres hatte, und der ebenfalls für die Störungen durch jeden der Planeten, Venus, Jupiter, Mars und Saturn, in Tafel II. angegeben ist. Das ganze Jahr hindurch bleibt also das zweite Argument dasselbe, und man thut am besten, aus den Tafeln die Gleichungen nicht für den Tag selbst zu nehmen, wofür man rechnet, sondern für den nächst vorhergehenden und nächst nachfolgenden, die sich unmittelbar vorfinden (wodurch also die Tafeln nur Tafeln zu einfachen Eingängen werden), und nachher zwischen die Aggregate aller vier Störungen zu interpoliren. In Beziehung auf diese trefflich angeordneten Störungstafeln sei es uns erlaubt, ein paar Kleinigkeiten zu bemerken. \textsc{Carlini} hat sie in der Voraussetzung berechnet, dass die mittlere Länge der Erde zu Anfange des Jahrs $= 100^\circ$ sei. Das heisst nun eigentlich so viel, man muss, \emph{wenn man es ganz strenge nehmen will,} in sie eingehen 1) mit der Anzahl von Tagen, die nicht seit dem Anfange des Jahrs, sondern seit dem Augenblick verflossen sind, wo die mittlere Länge der Sonne $= 100^\circ$ war; und 2) mit dem Werthe, den die mittlere Elongation des störenden Planeten, auch nicht zu Anfang des Jahrs, sondern in eben diesem Augenblick hatte. Der Unterschied ist allerdings nie irgend erheblich, \SourceLeaf{551}inzwischen hätte man doch wünschen können, es immer nach Gefallen in seiner Gewalt zu haben, die genau richtigen Argumente zu wählen; in Ansehung des ersten Arguments hätte dieses durch eine besondere Correctionstafel sich thun lassen, dergleichen wir zu unserm Gebrauch unserm Exemplare beigefügt haben, und das zweite Argument hätte eben so leicht für den genau richtigen Augenblick in Tafel II. angesetzt werden können, als für den Anfang des Jahres. Im Jahre 1804 ist z.\,B. die Correction des ersten Arguments 1.09, und die genauen Werthe von $E$, $F$, $G$, $H$, sind: \[ \begin{array}{r|l@{\;}l} E = 36.66 & \text{in der Tafel} & E = 36.55\\ F = 28.45 & & F = 28.34\\ G = 13.13 & & G = 13.09\\ H = 10.1 & & H = 10.1 \end{array} \]

Zweitens hätten wir die Tafel selbst noch um 10 Tage weiter ausgedehnt gewünscht, nemlich auf 370 Tage anstatt 360. Es wäre dadurch für den Fall, wo man für die letzten Tage des Decembers einen Sonnenort berechnen will, die (freilich nur kleine) Mühe erspart, die Argumente $E$, $F$, $G$, $H$, zugleich für das nächstfolgende Jahr mit auszuschreiben.

Wir haben, um unsern Lesern eine Idee von der grossen Bequemlichkeit dieser Tafeln zu geben, die Einrichtung derjenigen, welche zur Berechnung der Länge dienen, mit einiger Ausführlichkeit beschrieben. Der beschränkte Raum verstattet uns nicht, bei den übrigen Tafeln eben so umständlich zu Werke zu gehen. Wir bemerken also nur, dass die Tafeln zur unmittelbaren Berechnung des Logarithmen des \emph{Abstandes} der Sonne von der Erde, und die zur Berechnung der \emph{Breite} der Sonne, eben so geschmeidig angeordnet sind, wie die für die Länge. Ausserdem sind noch ein paar Tafeln für die Reduction der Ekliptik auf den Aequator, für die mittlere Rectascension der Sonne in Zeit, und zur Reduction der Sonnenörter vom mittleren Mittag auf den wahren, angehängt.

Die den Tafeln vorgesetzte Einleitung entwickelt ihre Einrichtung mit Klarheit und Eleganz. Aber einen vorzüglich hohen Werth erhält sie durch die Mittheilung vieler sinnreicher Kunstgriffe, um die Berechnung einer Ephemeride für die Sonnenörter auf ein ganzes Jahr möglichst zu erleichtern. In das Detail derselben einzugehen, erlaubt uns hier der Raum nicht: allein wir halten es für Pflicht, allen Astronomen, welche viel Sonnenörter zu berechnen haben, das Studium und die Benutzung dieser Vortheile angelegentlich zu empfehlen.
